\section*{ЗАКЛЮЧЕНИЕ}
\phantomsection
\addcontentsline{toc}{section}{Заключение}

В рамках курсовой работы выполнен численный анализ трибологических характеристик радиального подшипника скольжения с дискретным микрорельефом поверхности. В ходе исследования решены следующие задачи:

\begin{enumerate}
    \item Выполнен обзор литературных источников по теории гидродинамической смазки и методам лазерного текстурирования поверхностей трибосопряжений. Показано, что основным физическим механизмом повышения эффективности является микрогидродинамический эффект, возникающий на кромках углублений за счёт формирования локальных клиновых зазоров.

    \item Реализована методика численного решения уравнения Рейнольдса для подшипника конечной длины на основе метода конечных разностей с итерационной схемой Гаусса-Зейделя и верхней релаксацией. Расчёт проводился на сетке $500 \times 500$ узлов при критерии сходимости $\delta < 10^{-5}$.

    \item Определены поля давления и интегральные характеристики (нагрузочная способность $F$, коэффициент трения $\mu$, расход смазки $Q$) для гладкого подшипника и подшипника с углублениями в форме сферической шапки (тип~5) при эксцентриситетах $\varepsilon = 0{,}05 - 0{,}80$. Параметры подшипника: $R = 35$~мм, $c = 50$~мкм, $L = 56$~мм.

    \item Проведено сопоставление гладкого и текстурированного подшипников. При рабочем эксцентриситете $\varepsilon = 0{,}6$ установлены следующие изменения характеристик:
    \begin{itemize}
        \item нагрузочная способность возросла на $10{,}42$\%: с $F = 5028{,}5$~Н до $F = 5552{,}5$~Н;
        \item коэффициент трения уменьшился на $11{,}22$\%: с $\mu = 0{,}01538$ до $\mu = 0{,}01365$;
        \item расход смазки увеличился на $1{,}27$\%: с $Q = 0{,}2359$~л/с до $Q = 0{,}2389$~л/с.
    \end{itemize}
\end{enumerate}

Результаты исследования подтверждают, что нанесение углублений в форме сферической шапки с параметрами $a = b = 2{,}3$~мм, $h_p = 12$~мкм ($h_p/c = 0{,}2$) приводит к заметному улучшению трибологических характеристик подшипника. Совмещение прироста несущей способности со снижением трения при умеренном увеличении расхода смазки указывает на высокую эффективность рассматриваемого типа микрорельефа.

Профиль сферической шапки аппроксимирует сферический сегмент и характеризуется равномерным изменением кривизны по всей площади углубления. В отличие от эллипсоидального профиля, где кривизна различается вдоль главных осей, сферическая форма обеспечивает осесимметричное распределение локального давления, что способствует формированию стабильного микрогидродинамического эффекта без концентрации напряжений в отдельных направлениях.

На основании полученных результатов рекомендуется применение микрорельефа типа сферической шапки для подшипников скольжения нефтегазового оборудования, эксплуатируемого в режиме жидкостного трения. Наибольший положительный эффект достигается при повышенных эксцентриситетах ($\varepsilon > 0{,}5$), где прирост несущей способности и снижение потерь на трение проявляются наиболее существенно.
